\documentclass[a4paper,11pt]{article}

% Pacotes Principais -----------------------------------------------------------
\usepackage[portuges,brazil]{babel}
\usepackage[utf8]{inputenc}

% Formatação de capítulos ------------------------------------------------------
%\usepackage[Sonny]{fncychap}
%\usepackage{fncychap}
\usepackage{capitulos}

% Figuras e Imagens ------------------------------------------------------------
\usepackage{graphicx}
% Figuras lado a lado
\usepackage{epsfig}
\usepackage{subfigure}

% Utilizar H para inserir as imagens REALMENTE onde eu desejo
\usepackage{float}

% Fontes -----------------------------------------------------------------------
\usepackage[T1]{fontenc}
\usepackage{pslatex}

% Simbolos ---------------------------------------------------------------------
\usepackage{textcomp}
\usepackage{bbding}

% Tabelas ----------------------------------------------------------------------
%\usepackage{multicol}
\usepackage{multirow}
% Colorir a tabela
\usepackage{colortbl}
% Pacote hhline corrige os bugs das linhas que não aparecem com o colortbm
\usepackage{hhline}
% Tabelas com colunas de largura auto ajustável
\usepackage{tabularx}
% Notas de rodapé em tabelas (Pode-se usar o ambiente longtable também -
% Pesquisar exemplo com longtable)
\usepackage{threeparttable}
% Tabelas grandes
\usepackage{supertabular}

% Glossário --------------------------------------------------------------------
\usepackage[portuguese,noprefix]{nomencl}
\usepackage{makeglo}

% Outros pacotes ---------------------------------------------------------------
\usepackage{noitemsep}

% Comentários em bloco
\usepackage{verbatim}

% Hiperlinks
\usepackage{hyperref}

% Orientação de página
\usepackage{lscape}

% utilitários matemáticos
\usepackage{amsmath}
\usepackage{icomma}

% Evita o problema "too many unprocessed floats", colocando os campos
% 'flutuantes' em suas respectivas seções
\usepackage[section]{placeins}

% Referências ------------------------------------------------------------------
\usepackage[abbr]{harvard}	% As chamadas são sempre abreviadas
\harvardparenthesis{square}	% Colchetes nas chamadas
\harvardyearparenthesis{round}	% Parêntesis nos anos das referências
\renewcommand{\harvardand}{e}	% Substituir "&" por "e" nas referências

% Comandos gerais --------------------------------------------------------------
\newcommand{\titulo}{Sistemas de Tempo Real}
\newcommand{\autor}{Diogo Leite Rebouças}

% Configuração da fonte
%\renewcommand{\familydefault}{\sfdefault}

% Margens ----------------------------------------------------------------------
\setlength{\oddsidemargin}{3.5cm}
\setlength{\evensidemargin}{2.5cm}
\setlength{\textwidth}{15cm}
\addtolength{\oddsidemargin}{-1in}
\addtolength{\evensidemargin}{-1in}

\setlength{\topmargin}{2.0cm}
\setlength{\headheight}{1.0cm}
\setlength{\headsep}{1.0cm}
\setlength{\textheight}{22.7cm}
\setlength{\footskip}{1.0cm}
\addtolength{\topmargin}{-1in}

% Glossário --------------------------------------------------------------------
\makeglossary

% Capítulos --------------------------------------------------------------------
% Não aparecer o número na primeira página dos capítulos
\newcommand{\mychapter}[1]{\chapter{#1}\thispagestyle{empty}}
\newcommand{\mychapterstar}[1]{\chapter*{#1}\thispagestyle{empty}}

% Comandos matemáticos ---------------------------------------------------------
% Implicação em fórmulas
\newcommand{\implica}{\quad\Rightarrow\quad} %Meio de linha
\newcommand{\implicafim}{\quad\Rightarrow}   %Fim de linha
\newcommand{\tende}{\rightarrow}

% Fração com parenteses
\newcommand{\pfrac}[2]{\parent{\frac{#1}{#2}}}

% Transformada de Laplace e transformada Z
\newcommand{\lapl}{\pounds}
\newcommand{\transfz}{\mathcal{Z}}

% Sequências
\newcommand{\sequencia}[4]{$#1_{#2}$, $#1_{#3}$, \ldots, $#1_{#4}$}

% Outros ----------------------------------------------------------------------
\newcommand{\chave}[1]{\left\{#1\right\}}
\newcommand{\colchete}[1]{\left[#1\right]}
\newcommand{\parent}[1]{\left(#1\right)}

\newcommand{\rhoagua}{\rho_{\tiny \text{\tiny H}_2\text{\tiny O}}}

\let\D\displaystyle
\newcommand{\reg}{\textsuperscript{\textregistered}}
\newcommand{\soft}{\textit{software}}
\newcommand{\Soft}{\textit{Software}}

% Imagens exportadas pelo gnuplot ----------------------------------------------
\graphicspath{{imgs/eps/}}


\pagestyle{empty}
\setlength{\parindent}{12mm}
%\usepackage{indentfirst}

\begin{document}

\begin{table}[H]
\centering
\begin{tabular}{cl}
% Primeira Linha
\multirow{7}{*}
{
\includegraphics[width=1.65cm]{imgs/uern}
} & \\
& Governo do Estado do Rio Grande do Norte\\
% Segunda Linha
& Secretaria de Estado da Educação, da Cultural e dos Desportos -- SECD\\
% Terceira Linha
& {\sc Universidade do Estado do Rio Grande do Norte -- UERN}\\
% Quarta Linha
& Pró-Reitoria de Ensino e Graduação -- PROEG\\
% Quinta Linha
& Ciências da Computação -- Trabalho teórico para 3\textordfeminine\ unidade (10
pontos)\\
& 
\end{tabular}
\end{table}

\begin{center}
\Large
\uline{{\sc Trabalho Teórico de Pesquisa}}
\end{center}

\section*{Informações gerais}

Apesar de sabermos que o campo de Sistemas Multimídia abrange diversas áreas de
atuação, não conhecemos em detalhes que áreas são essas. Por ser uma matéria que
continuamente evolui, a cada dia surgem mais segmentos de pesquisa. Pesquisar
sobre esses segmentos pode abrir nossa mente para várias coisas e,
principalmente, fazer com que tenhamos ideia\$ inovadora\$.

Assim sendo, nesta terceira unidade é proposto um trabalho de pesquisa em que
vocês deverão trazer a tona os temas mais recentes da área de sistemas
multimídia. Para isso, deverão {\bf elaborar um} texto no formato de {\bf
artigo e uma apresentação}.

O trabalho (artigo + apresentação) deverá ser entregue {\bf até às 23:59h da
data estabelecida} em sala de aula, podendo ser feito maneira individual ou em
grupo de no máximo 4 (quatro) integrantes. Após estabelecida a data, \uline{{\bf
não haverá prorrogação do prazo}}.

A avaliação do trabalho será realizada em dois dias e todos os grupos deverão
estar presentes. A participação dos grupos, dando sugestões ou tirando dúvidas
sobre o desenvolvimento, também será levada em consideração na avaliação.

A nota será a mesma para todos os componentes do grupo, havendo penalização para
cada dia de atraso na entrega. Caso algum componente do grupo não possa
comparecer no dia da apresentação, a falta deverá ser justificada junto à
coordenação, de acordo com as normas da UERN. Nos casos em que a {\bf falta não}
for {\bf justificada}, o aluno receberá {\bf nota 0.0 (zero)}.

\section*{Os critérios de avaliação}
Com o intuito de incentivar a produção de documentos em \LaTeX, será fornecido
um {\it artigo modelo} para que vocês modifiquem o conteúdo e adequem ao tópico
selecionado. Além do modelo citado, serão fornecidos também os arquivos da
apostila e da apresentação do minicurso de \LaTeX. 

Dessa forma, não será permitida a entrega dos artigos produzidos através do
Word, BrOffice, OpenOffice ou qualquer outra ferramenta do gênero. Os trabalhos
entregues nesse formato perderão de antemão 30\% da nota. Qualquer dúvida não
esclarecida com o material fornecido pode ser retirada por e-mail ou
pessoalmente.

As apresentações, por outro lado, poderão ser desenvolvidas em qualquer tipo de
ferramenta. Contudo, as apresentações desenvolvidas em \LaTeX\ receberão um
bônus de 10\% sobre a nota final. Logo, os seguintes critérios serão levados em
consideração para composição da nota:

\begin{enumerate}
    \item Desenvolvimento do artigo em \LaTeX (30\% da nota)
    \item Apresentação do tema (55\% da nota)
    \item Participação em outras apresentações (15\% da nota)
    \item Desenvolvimento da apresentação em \LaTeX (10\% extra)
\end{enumerate}

Os arquivos referentes ao artigo e a apresentação deverão ser organizados em uma
estrutura de pastas, compactados e enviados para {\tt diogolr@gmail.com} com o
Assunto {\tt [SistMult-Natal] Trabalho - Grupo X} para a turma do Campus de
Natal, ou {\tt [SistMult-NC] Trabalho - Grupo X} para a turma do Núcleo de Nova
Cruz.

\section*{Nota final}
A nota final será composta pela soma dos critérios acima estabelecidos e pela
penalização por dia de atraso, conforme Eq. \ref{eq:nota}.

\begin{equation}\label{eq:nota}
N_f = \D 1,25^{-D_a} \times \sum_{i = 1}^3 C_i
\end{equation}

\noindent em que $D_a$ se refere aos dias de atraso contados a partir da data
estabelecida e $C_i$ corresponde ao $i$-ésimo critério.

\section*{Tópicos para os trabalhos}

Os seguintes tópicos são sugeridos para a elaboração do trabalho:

\begin{itemize}
    \item 3DTV
    \item Linguagem SMIL
    \item Ferramentas para o desenvolvimento de aplicativos para TV Digital
    \item H.264/Mpeg-4 SVC – Scalable Video Coding
    \item Padrões de comunicação em TV Digital (brasileiro $\times$
          estrangeiros)
    \item Medição de qualidade de vídeo em redes de pacotes (QoE -- {\it Quality
          of Experience}, modelos para qualidade de vídeo)
    \item Redes de Distribuição de Conteúdo (CDN)
    \item Vídeo por demanda
    \item Mídias de armazenamento
\end{itemize}

{\bf Os alunos não estão restritos à esses tópicos}. Qualquer outro tópico,
desde que relacionado com o processamento de imagens, áudio ou vídeo, poderá ser
enviado por e-mail até 20 dias antes da data da apresentação, estando sujeito à
aprovação pelo professor.

Após a escolha do tema pelo grupo, qualquer modificação só poderá ser realizada
até 20 dias antes da data da apresentação (através de e-mail).

\end{document}
